\chapter{Witness, Reality, and Decision}

Chapter 4 closed with an apparatus that could keep an account answerable: it could receive an origin, carry a receipt, gather a contribution, set a scale, and hold overhead as a conserved pressure rather than waste. What that apparatus could not yet say is what makes such an account true. The distinction is sharp, and the whole chapter turns on it. An answerable account is one that can be challenged; a true account is one that survives the challenge, and survives it for the right reason rather than by accident, fatigue, or force of repetition. This chapter is about that reason. It carries truth through the apparatus one demand at a time --- as a witnessed receipt rather than an owned possession, as a variation tested where the claim stands, as a reading carried along paths without losing what earned it, as a calibrated reading that stops at a finite place, and finally as a report that closes on the claim it began with. It is the chapter in which the apparatus, having learned to measure, learns what it takes for a measurement to be true.

None of these demands is imported from outside the apparatus. Each is forced by the same receipt the earlier chapters earned, and each asks of that receipt a question the previous demand left open. Witness asks whether the account can be received by a position other than the one that asserted it. Locality asks whether it can be tested where it stands. Mediation asks whether the test survives being carried elsewhere. Calibration and stopping ask whether the carried reading can end in a value rather than run forever. Closure asks whether the ended reading answers the claim that set the whole passage going. The chapter is the working-through of these five questions in the only order in which they can be asked, because each becomes askable only once the one before it has an answer.

The chapter's title names three things, and the five demands distribute among them. Witness is truth received rather than owned, the first demand and the condition for every later one. Reality is the claim meeting more than the apparatus's own order: a local variation answers to something the apparatus does not control, and a mediated path carries the claim into positions that can resist it, so the second and third demands are where reality presses on the account. Decision is the apparatus committing rather than drifting: a calibrated stop decides that the reading has gone far enough to read, and an original closure decides that the report has answered the claim it began with, so the fourth and fifth demands are where the apparatus must choose. Witness opens the chapter, reality tests it through local variation and travel, and decision closes it at a finite, answerable place.

The first demand refuses ownership of truth. A claim is not true because the apparatus holds it; it is true because a witness can receive it and answer for it. The asserting position cannot be its own witness, for a claim that only its author can confirm has never left the place where it was made. So the first section treats scientific truth as a witnessed receipt rather than a possession --- something the apparatus tanges into a record that a later, independent position can read, challenge, and return to, not something owned by whoever first pronounced it. The earlier bridge had already prepared the route by which a witness might challenge the apparatus; witness is the encounter that actually walks that route. Witness in this sense is not a second opinion added to a finished truth; it is part of what makes the claim a truth at all, since the relation between the position that asserts and the position that receives is where public standing lives, and neither position holds that standing alone. Truth that can only be held is not yet public; truth that has been witnessed has entered a record where another position can press it, and a claim that cannot bear that pressure was never the kind of thing the apparatus had a right to call true.

The second demand is locality. A witnessed claim must be testable where it stands, because a finite apparatus cannot check a claim everywhere at once; it can only read how the claim responds to a small, deliberate variation in what it depends on. So the second section funges a nearby perturbation through the claim and reads what comes back, refining what counts as a response from a crude difference toward a genuine local rate --- the Newton-to-Gateaux-to-Frechet progression that sharpens the question of what a claim does under a controlled nudge. The refinement is a gain in fidelity, not a flight into the infinite: the apparatus asks only for a response it can actually exhibit, not for a limit it could never finish taking. Truth is tested locally before it is trusted globally, and the local variation is where a claim first answers for itself rather than for the reputation of the position that made it.

The third demand is mediation. A claim tested in one place is worth little if it cannot reach another, but reaching another by simply asserting it again would leave the new position with a verdict and no way to check it. So the third section follows truth along paths, through the universal mediation by which a local reading travels carried rather than re-asserted. A path funges the local reading onward together with the terms that earned it, so that a distant position inherits not the conclusion alone but the route back to the receipt that supports it. A truth that cannot travel is parochial, confined to the position where it was first read; a truth that travels by mediation keeps its receipts as it goes, and remains as answerable at the end of the path as it was at the start.

The fourth demand is calibration and stopping. A reading that never stops is not a measurement but a deferral, and an account that defers forever answers nothing. So the fourth section gives the apparatus a calibrated clock, a stop, and a measured output: it calibrates the reading against a fixed floor that says when a difference is small enough to ignore, tanges the running process at that floor, and funges out a finite value that another position can receive. The stop is not an interruption of the measurement but its completion --- the moment at which an ongoing process becomes a result. Without calibration the apparatus could not say when it had read enough; without the stop it could never hand anything on. Measurement requires an end, and the end must be earned at a finite place rather than promised at an unreachable one.

The fifth demand is closure, and meeting it completes the apparatus. A report has closed originally when it returns to the claim that began it and answers that claim from the terms carried in between, rather than from its own length or the assent of its audience. So the fifth section is that original closure: the apparatus, having witnessed a claim, tested it locally, carried it by mediation, and stopped its reading at a calibrated place, returns to the opening claim and closes on it, inferring its standing from its own earned record. Closure on the original claim is what distinguishes a measurement that answers from a process that merely halts; the report ends where it began, and ends there having earned every term it needed to get back.

Together these five demands carry truth from witness to inference without ever letting the apparatus claim to own it. The chapter does not assert that the apparatus possesses the truth; it shows the apparatus witnessing a claim, testing it where it stands, carrying it without loss, calibrating and stopping its reading, and closing on the claim it set out to make. Each demand answers to the same receipt, and the last returns that receipt to the first. By the chapter's end the apparatus has closed on its own truth --- not as a possession it holds against challenge, but as a witnessed, carried, stopped report that answers, originally, the claim it began with. What the apparatus could not do at the close of Chapter 4, it learns here to do: not merely to remain answerable, but to be answered.

\section{Scientific Truth as Witness, Not Ownership}

Chapter 4 ended with an apparatus that could keep an account answerable. It could receive an origin, carry a receipt, gather a contribution, set a scale, bear a load, and preserve the overhead that lets a later hand return to the same scene. Such an apparatus matters because it refuses a loose claim: it does not let a statement float away from the route that first gave it public shape, and it keeps the mark beside the carried thing, the comparison beside the standard, and the burden beside the support that accepted it. Yet this whole order still stops short of truth. The apparatus can answer for how a claim passed through it, but answering for the passage does not by itself decide what the claim deserves when another position comes to test it. A receipt can say where an account began and how it was preserved; the receipt cannot become the encounter that faces the account as a claim. Chapter 5 begins at that threshold. The apparatus has learned to keep a claim returnable; now it must ask what sort of witness can meet the return without taking possession of the truth it seeks.

Scientific truth enters here as a discipline of witness, not as a form of ownership. No observer owns a truth because the observer first noticed a mark. No bearer owns it because the bearer carries the receipt. No group owns it because many voices repeat the same account. No office owns it because it keeps a seal or a public name. Ownership can protect a vessel, guard a record, or assign responsibility for a carried account, but ownership cannot turn the account into truth. The reason lies in the relation. A claim asserted at a position and held there is a closed loop,
\[ \text{ownership:}\quad P \xrightarrow{\;c\;} P \qquad (\text{asserter} = \text{receiver}), \]
in which the position that makes the claim is the only one that receives it. In pigeonhole terms, the asserter and witness are two roles that cannot occupy one hole; the witness needs its own receiving position, so finitude forces the route from $P$ toward a distinct $Q$. Private certainty fails for exactly this reason: a person may feel no doubt and still leave the claim with no route another position can challenge. Adoption fails the same way when it asks agreement to do the work of witness, since a crowd can copy a phrase or inherit a practice while the claim itself slips away from the receipt that should answer for it. Scientific truth therefore has a positive shape: a claim enters a witnessed route, offering itself with enough carried order that another participant can return to the mark, inspect the receipt, raise a challenge, refuse a broken route, and preserve a confirmation only when the route survives.

Witness makes the route public, and publicity brings danger with it. A witness does not look at a claim from a safe distance; a witness accepts an obligation to meet the claim where the apparatus says the claim can answer, and tanges it there --- asking what entered, which receipt carried it, what comparison faced it, which refusal the scene allowed, and what remains after the challenge. The claim cannot demand trust before this meeting. It cannot wear the witness's name as decoration, and it cannot treat attention as endorsement. So the witness must resist two failures at once: it must not let the claim hide in private conviction, and it must not let public attention harden into a false maker of truth. A crowd can gather and still leave the claim untested. A confident voice can speak well and still fail to keep the comparison available. A standing practice can remember that a comparison once held while losing the comparison itself. Witness answers these dangers by returning to the carried scene and asking for the receipt, the rule, the visible challenge, and the place where refusal can still matter.

This public obligation also explains why witness cannot become borrowed authority. A claim may travel with a strong account --- one that names a careful origin, a long lineage of repeated practice, or a route that many hands have used. Those facts can help a witness find the route, but they cannot replace it. Borrowed authority asks the witness to honor the reputation instead of testing the carried claim: to bow before continuity, to accept a receipt because the receipt wears an old name, or to treat repetition as if repetition had already done the work of challenge. Scientific truth refuses that shortcut. It welcomes memory only when memory leads back to what the apparatus keeps available, welcomes standing practice only when practice keeps the claim open to refusal, and welcomes public confidence only when confidence stays accountable to the scene another witness can enter. Truth does not grow out of borrowed status; it travels through a route that keeps status answerable to marks, receipts, challenges, and possible confirmation.

Witnessed receipt names the heart of that route. The receipt must remain more than a souvenir of arrival: it must show what claim entered the apparatus, how the claim moved into public reach, what route lets another participant return, and what sort of challenge the claim can meet. A receipt that names only an origin leaves the later witness too little; a receipt that names only a result cuts the route away from the mark that made the result answerable; a receipt that keeps the route but allows no refusal funges witness into ceremony. For the receipt to do its work it must stay open on several faces at once. It must expose the claim without pretending exposure settles it, preserve the route without treating the route as a possession, let challenge press against the account without making every refusal final, and let possible confirmation remain possible without turning confirmation into a private treasure. The receipt witnesses only because it holds the claim, the route, the challenge, the refusal, and the possible confirmation in one returnable order.

Refusal gives that order its bite. If every challenge had to end in acceptance, witness would merely decorate the claim with public ceremony; if every challenge had to end in rejection, the apparatus would confuse suspicion with discipline. Refusal matters because it keeps the route honest at the point where the claim could fail. The witness may find a missing receipt, a broken route, an unclear comparison, or a claim that asks the apparatus to carry more than the scene can answer for, and in that moment refusal protects truth from eagerness. It also protects the claim from rumor, because a refusal names the place where the route lost its answer rather than dissolving the whole account into distrust. A later participant can return to the same wound in the route --- repair the missing receipt, repeat the comparison, or show why the earlier refusal pressed on the wrong point. Refusal therefore belongs inside witness, not outside it: it keeps the apparatus public enough to fail and precise enough to try again.

The claim itself must remain visible within that order. If the apparatus keeps only marks and no claim, a later witness cannot tell what faced the challenge; if it keeps only a claim and no route, the claim asks for belief without answerability; if it keeps only a route and no refusal, the route becomes a procession that never risks meeting an objection. Scientific discipline requires all of these together. The witness should be able to ask what the claim said, where it entered, which receipt carried it, what challenged it, what refusal the challenge made possible, and what remained when the apparatus returned to the scene. None of these questions owns truth; each protects the difference between an account that merely circulates and a claim that survives witness. The apparatus needs no lord of truth. It needs a route where the claim can lose, persist, or return with a confirmation that another witness can still examine.

This is why the discipline cannot stay with the observer. The observer may begin the route by noticing a mark, but the observer's attention does not give the mark final standing; attention begins a responsibility, not a verdict. The bearer cannot keep truth inside the carried receipt either: the bearer can preserve the receipt and move it from one scene to another, but carriage alone does not test the claim. Nor can an audience complete the route by receiving the account; listening, repeating, applauding, or doubting cannot substitute for a returnable encounter. What the discipline requires instead is a directed passage from the position that asserts the claim to a distinct position that answers for it,
\[ \text{witness:}\quad P \xrightarrow{\;c\;} Q \xrightarrow{\;\text{answers for } c\;} \text{record}, \qquad Q \neq P, \]
in which the route funges the claim from its asserting position $P$ to a receiving position $Q$ that can challenge, refuse, or confirm it, and whose answer enters the shared record. The single inequality $Q \neq P$ carries the whole section: a position cannot witness its own claim, and a truth that never leaves $P$ has never been tested at all.

Return also gives even the first witness a limit. The first witness may perform the encounter carefully, but the encounter does not become final because it came first; a later witness must be able to ask the scene to answer again. The apparatus must therefore hold more than testimony about a past success --- it must hold the shape of the encounter: the claim that entered, the receipt that carried it, the challenge that pressed it, the refusal that could have stopped it, and the confirmation that travels only as long as those pieces remain available. This makes the first witness accountable to later witnesses without making later witnesses superior by number. Sequence helps only when each return keeps the route intact.

The returnable arrangement gives scientific truth its severity. It does not make truth cold or ownerless by neglect; it makes truth ownerless by refusing to let any single holder close the route. The apparatus protects the mark from private possession, the receipt from mere ceremony, the challenge from collective applause, and the confirmation from becoming a locked trophy. A claim that cannot return may still inspire speech, habit, or confidence, but it has not entered this discipline. A claim that can return must expose itself to a witness who may ask for the route again: it must let the receipt stand beside the challenge, let refusal remain visible even when the claim survives, and let confirmation travel as a public answer rather than a prize held by the first observer. This is the kind of truth the chapter now follows --- truth carried through witness, never owned by the carriers who help it travel.

Once truth needs witness, the apparatus faces a sharper question. A witnessed route cannot hang in empty speech; it must meet what the claim concerns and ask how the surrounding order constrains the return. A witness can demand the receipt and challenge the route, but the route still has to answer from somewhere --- it has to show how the claim touches a world that resists fantasy, how a place of return limits what the claim may say, and how a possible change can test the claim without dissolving the account. Those questions belong to the next section. This one closes with the first rule of Chapter 5: truth belongs to no owner, observer, bearer, crowd, or institution. A claim is true not because some position holds it, but because a different position can receive it and answer for it; truth is what survives the distance between the two.

\section{Variation and Locality}

The witnessed route from the previous section gives truth a public discipline, but it still leaves a question unanswered. A witness can ask for the claim, demand the receipt, press the route, and keep refusal available, yet none of that tells the apparatus where the claim must answer. A claim that floats in speech can survive any challenge by moving the place of answer whenever pressure arrives --- saying that the witness misunderstood the scene, that the receipt belonged somewhere else, or that the refusal missed the real point. The apparatus cannot allow that escape. Once a claim enters witness, the route needs a bounded place of return: the witness must be able to bring the receipt back to a scene where the claim either meets resistance or loses its authority to travel.

Reality names that resistant answerability. It does not name a possession, a private vision, a single authority, or a larger voice that settles every dispute from above. Reality enters the apparatus as the surrounding order that can say no to a claim, and that order need not explain itself in full before it can refuse; it need only give the claim a place where the claim's own carried terms face something that does not bend for convenience. The claim returns to the scene, the receipt returns to the mark it names, and the surrounding order can withhold the answer the claim wanted. That withholding is the point. If reality merely echoed the claim, witness would add ceremony without danger; if reality merely crushed every claim, witness would turn discipline into denial. Resistant answerability keeps both failures away, letting a claim meet an order that can answer, fail to answer, or answer only after the route narrows its demand.

This resistance also prevents reality from becoming another owner of truth. An owner closes a route by claiming the account as a possession; resistant answerability keeps the route open by demanding return, asking the claim to come back to the scene and show what it can carry. The apparatus therefore treats reality neither as a hidden master nor as a decorative background, but as the pressure that makes return cost something. A claim that concerns a measurement must meet that measurement where the receipt can bear the weight; a claim that concerns a count must meet the count where another hand can repeat it; a claim that concerns a refusal must meet the refusal where the route broke or survived. The surrounding order does not replace witness; it gives witness a place to stand.

Locality names that place of standing. A local place is not a tiny world sealed off from every other scene; it is the bounded return-place where the apparatus can bring the claim, receipt, challenge, and refusal together without losing which part of the order answers. Locality protects the witness from empty generality. A claim cannot answer from nowhere in particular and still demand the privileges of return; it must say enough for a witness to find the place where the receipt can press. The local scene gathers the relevant mark, the carried route, the available refusal, and the possible answer --- and it gathers nothing else. Its power comes from that restraint.

That restraint lets locality do two things at once. It limits what a claim may say: a claim that answered in one scene cannot pretend that every scene has answered, and a receipt that survived one challenge cannot pretend that every future challenge has lost its bite. And it lets an answer travel without becoming shapeless, since a receipt that keeps the local conditions visible lets a later witness ask whether another scene carries the same relevant order. The apparatus never leaps from one successful return to a total claim; it carries the local answer as a disciplined piece of evidence --- this claim met this resistance, under these terms, with this refusal still visible. Without a local place a claim inflates after every survival, turning a local success into a universal one, and the named local limit is exactly what keeps a strong answer from swelling into a loose proclamation. The limit does not weaken the answer; it makes the answer honest enough to travel.

Locality also gives refusal a sharper form. In the previous section refusal protected truth from eagerness; here it gains a place. A local refusal does not merely report that the witness disbelieved the claim --- it says where the scene withheld answer: the mark failed to line up with the receipt, the route asked for more support than the scene could give, the comparison could not return to the same standard, or the claim changed its demand when the witness brought it back. Each such failure leaves a local wound that another participant can inspect, repair, or show to have pressed the wrong feature. Locality keeps a failure from spreading into mere distrust and keeps a survival from swelling into empty confidence.

Once locality holds the claim in place, variation can begin. Variation does not mean that anything may change in any direction; it means the apparatus can alter one pressure while keeping enough of the route fixed for the local scene to answer. The apparatus tanges the claim at the bounded scene and reads its response to a perturbation of definite size,
\[ \Delta_h\, c(x) \;=\; \frac{c(x+h) - c(x)}{h}, \qquad h \neq 0, \]
the change in the claim's answer when one input is moved by a finite step $h$. This reading is something the apparatus can exhibit at a chosen $h$; it never has to drive $h$ toward a limit it could not finish taking. The simplest variation asks exactly this: if the claim shifts here, does the scene still answer? A receipt may stay the same while a different burden presses against it, or a route may keep its origin while a later hand changes the demand. This first pressure separates a claim that merely survived one performance from a claim that can face a nearby change without losing its route.

The local scene does that work --- not the witness by preference, and not the claim by its own report. The apparatus funges the changed demand back to the bounded scene and asks the surrounding order to answer. Sometimes the scene accepts the change because the receipt still reaches the same mark; sometimes it refuses because the changed demand breaks the route; sometimes it answers only after the claim gives up a larger promise and keeps a smaller one. In each case variation reveals what bare repetition cannot. Repetition asks whether the same route can return; variation asks how much pressure the route can bear while it still remains the same route.

Stronger variation adds direction. A witness can ask not only whether a single change survives, but how the claim responds when the apparatus funges a perturbation along a named line of pressure,
\[ D^{h}_{v}\, c(x) \;=\; \frac{c(x + h v) - c(x)}{h}, \]
the response read along a direction $v$ at the same finite step $h$. Direction matters because not every change threatens the same part of the route: one line presses the claim's origin, another its receipt, another the refusal the earlier scene allowed. A directed reading keeps those pressures distinct and tells the witness which part of the local order carried the answer and which part still waits for another return. Classical analysis names this refinement after Newton, Gateaux, and Frechet --- the plain difference, the response along a chosen direction, and the response that agrees across every direction --- but the apparatus keeps each as a finite reading it can exhibit, taking the gain in fidelity without the passage to a limit those names classically complete. Both $\Delta_h$ and $D^h_v$ are readings at a definite resolution, named more sharply, not limits pursued to vanishing.

A directed reading also introduces residue. A claim may answer along one line and still leave something unsettled beside it: the receipt may survive a changed burden while the comparison that gave the burden meaning has shifted, or the comparison may survive while the refusal moves to a different edge of the scene. This leftover does not embarrass the apparatus; it is the honest sign that the route answered only the pressure it actually faced. Residue keeps a local success from pretending to close every side of the claim, and it gives the next witness something precise to carry forward --- not a vague doubt, but a named remainder of the local encounter.

The strongest discipline in this section asks for compatibility among pressures. A claim that survives one change and one directed reading still has to show how those answers sit together. The apparatus brings two directed readings to the same bounded scene and asks whether the route keeps one account rather than splitting into convenient pieces. This is the last gain in fidelity the section reaches: not a single direction but the demand that the directed readings agree, so that the local response does not depend on which line of pressure the witness probes first. Compatibility does not erase residue; it asks whether the residues themselves can share a returnable order. If one reading says the claim survives and another says the route breaks, the witness must learn whether those results concern different features, competing demands, or an actual contradiction in the carried account. Local compatibility makes that question public.

This gives variation its proper modesty. It opens no unlimited theory of change and imports no finished analysis from a larger discipline; it works inside the apparatus already earned by witness, receipt, refusal, and bounded return. A local scene receives a pressure, answers or refuses, leaves a remainder, and asks whether several pressures can share one accountable route --- and that is enough for the chapter here, which does not yet need travel through many scenes but only a claim that faces one scene honestly. The same modesty preserves the witness. A witness who changes every condition at once no longer knows which return the receipt answered; a witness who refuses every change learns only that the first encounter can repeat. The discipline lies in the middle: change one pressure, keep the route recognizable, name the residue, and ask whether a stronger pressure still shares the account. That grammar lets the witness say what changed, what stayed available, what resisted, what remained, and what another hand may try next, while no single holder owns the answer.

The section therefore closes with a narrower but stronger demand than the one it opened. Truth cannot stop at witness, because witness without place can still drift: reality gives the route resistant answerability, locality gives that resistance a bounded return-place, variation lets the witness press the claim without dissolving the account, residue records what the pressure did not settle, and compatibility asks whether several local pressures can still keep one account. The next problem follows from this one: if a claim can answer in a local place, how can a route carry that answer to another place without losing the terms that made the first answer honest? That question belongs to the next section. A claim earns the right to be trusted somewhere only by first showing exactly where it can be made to answer; locality is the discipline that turns a claim's survival into a fact about a place rather than a mood about the claim.

\section{Paths and Universal Mediation}

The local return from the previous section gives a claim a place where it can answer. That place protects the witness from empty speech, but it raises a new question. A receipt that answers in one scene may need to travel to another, and a later participant may not stand in the same place, hold the same instrument, or face the same refusal. Still the claim may ask for passage. If it travels by leaving locality behind, it loses the very discipline that made its first answer honest. The route must move without pretending that motion grants a view from everywhere.

A path names that carried route between local scenes. It does not name a track across hidden territory; it names the way a witness keeps a claim answerable while moving it from one bounded return-place to another. The receipt does not travel as a loose slogan. It travels with the mark that called it forth, the refusal that gave it cost, the local limit that kept it honest, and the question that another scene may answer in its own way. A path therefore preserves difference as much as passage: the second scene need not repeat the first, but it must receive enough of the route to know what claim asks for answer, what receipt travels with it, and what refusal still has authority.

This makes a path stricter than a record of arrival. A bare report can say that a claim went from here to there; a path must show how the claim kept its accountable terms during that passage --- which part of the earlier answer belonged to the claim, which part belonged to the local scene, and which part remained unresolved. Without that memory, travel turns survival into exaggeration, and the claim treats one successful return as permission to speak everywhere. A path resists that inflation by carrying the local limit as part of the route, so a later answer can strengthen the claim only by meeting its own place of return.

Mediation names the discipline of that carriage. It does not erase the distance between scenes; it gives the distance a public rule of passage. A path is a finite chain of positions, and the mediation along it composes the local carries that funge the receipt from one position to the next,
\[ m_\gamma \;=\; m_n \circ \cdots \circ m_1 \;:\; P_0 \longrightarrow P_n, \qquad n < \infty, \]
where each $m_i$ preserves the conditions that gave the receipt its weight. The composition is finite --- a path is a sequence of definite carries, not a continuous transport through a space the apparatus could never traverse step by step --- and the receipt is what travels: $m_\gamma$ carries the route from $P_0$ to $P_n$ without re-asserting the verdict, so the later scene answers on its own terms while the earlier answer stays traceable. One place answers, another receives, and the claim keeps enough continuity for both refusals to remain intelligible.

The discipline matters most when a claim crosses unlike scenes. A claim may leave one bounded place and enter another, leave that one and meet a third, leave one witness and face another, and each scene can ask a different question and refuse in a different way. Mediation does not flatten those differences; it tanges the carried receipt at each new scene and asks whether it still gives that scene something determinate to answer. If the later scene cannot tell what survived, mediation fails; if the later scene must obey the first without its own answer, mediation also fails. The middle work succeeds only when passage and local answerability hold one another in tension.

Universal mediation names the compatibility discipline that belongs to any admissible passage of this kind. The word universal does not mean a rule above every scene, a sight from everywhere, or an owner above every place of return. It means that wherever the apparatus permits a route to travel, the passage must honor the same public demand: keep the claim's terms visible enough for each local scene to answer or refuse. A route cannot claim special privilege because it prefers one scene to another, cannot discard a refusal when the refusal becomes inconvenient, and cannot draw authority from distance. Universal mediation keeps the demand steady across passage without turning passage into a rule from nowhere.

That steadiness lets different witnesses compare routes without pretending that they share one final place. One may remember the first return, another inspect the second, a third ask whether the receipt survived the transfer; their work joins because mediation gives them a common demand --- show what traveled, what changed, which local answers still count, and which refusals remain open. Compatibility grows from that demand: it requires not identical scenes but enough accountable carriage that a later participant can pick up the route without swallowing the whole origin of the claim. The route thereby gathers memory. Each successful passage funges a trace into the shared record --- what the claim survived, where it narrowed, where the surrounding order withheld answer --- so the claim gains weight not because everyone agrees with it, but because its passage exposes a record that later hands can inspect, continue, or refuse.

That accumulated route memory closes the short hinge of this section. Locality gave the claim a place to answer; a path lets the answer travel while carrying the local limit; mediation keeps each later scene able to answer without surrendering its own terms; universal mediation names the compatibility every admissible passage must honor. The next question no longer asks whether passage can occur, but what arrests the route, what fixes a baseline for its report, and what lets the apparatus say the route has returned enough for public use. A truth that cannot travel is parochial, and a truth that travels without carrying its receipts is rumor; mediation is the narrow discipline that lets a claim move and stay answerable at once.

\section{Calibration, Halting, and Measured Output}

The previous section left a route at the edge of public use. A witnessed claim could travel from one local scene to another while carrying receipt, refusal, and limit, but that passage now needs a stricter end. A route that never stops cannot answer; it only continues. A route that stops without a baseline cannot show what its answer answers to. A route that reports without bounded terms cannot tell later hands what survived, what narrowed, and what still waits for return. The apparatus therefore needs three linked disciplines: it must fix a baseline for comparison, halt the extension of the route at an answerable place, and emit a bounded report that later witnesses can inspect, dispute, and repair without reopening the whole passage.

Calibration names the first discipline. It gives the route a public baseline, not a hidden authority, and it does not decide the claim by itself --- it gives the witness a stable place from which two route-states can meet. Before calibration, a later scene may see that something changed without knowing whether the change belongs to the claim, the scene, the witness, or the route's burden. Calibration narrows that confusion by fixing a floor: a tolerance $\varepsilon$ below which a difference between route-states no longer counts as disagreement. Two readings the route carries, $r_m$ and $r_n$, meet under calibration when
\[ d(r_m, r_n) \;\le\; \varepsilon, \]
with $d$ the difference the apparatus can actually measure. The floor is fixed in advance and held steady, so a later hand can check the same comparison the witness made; the baseline earns its authority because the apparatus can return to it, not because it stands above the route.

This public baseline also protects locality. A baseline that no later witness could revisit would become another form of private ownership; a baseline that absorbed every local difference would become another false view from everywhere. Calibration must do neither. It lets the later scene answer on its own terms while keeping enough common footing for the earlier receipt to stay legible. The apparatus may carry a claim through several local places, but each place still needs a way to say what changed against a shared point of return, and calibration supplies that point without making it mysterious. It does not turn the witness into an oracle; it gives the witness a floor another hand can check.

With the floor fixed, comparison becomes accountable. The apparatus tanges each route-state against the calibrated floor and asks what it can still carry. The comparison here is not a theory of all valid reasoning but the local ordering of route-states under a shared demand: one state may carry less burden than another, one may preserve the receipt while another loses the refusal that made the receipt honest, one may answer a narrower claim while another overreaches. Accountable comparison asks which state still answers the route and which has drifted from the terms of return. It needs no grand law of thought, only a public way to say that this route-state still stands before that one under the same carried question.

This matters because travel multiplies appearances. A route may look stronger merely because it crossed more scenes, weaker because a later scene exposed a limit the first did not test, or different because the burden changed while the words stayed the same. Calibrated comparison stops those appearances from passing as verdicts: it weighs each route-state against the same floor and asks what each can still carry, and the answer may favor the later state, restore the earlier one, or mark a residue between them. It also gives refusal a fair address, keeping a refusal that presses the starting mark distinct from one that presses the carried receipt, the size of the burden, or the way the route changed between scenes. Without that ordering the apparatus could confuse a louder report with a stronger one, treat a later refusal as having undone every earlier answer, or treat an earlier answer as having silenced every later refusal. Accountable comparison prevents those mistakes by making the route-state itself answerable.

Halting names the second discipline. The route halts when the apparatus stops extending it and exposes what can count as answer at that point. This promises no solution for every possible procedure and no guarantee that every search must end; it says something more local --- this route, under this floor and this comparison, has reached a place where further extension would hide the answer rather than clarify it. Concretely, the apparatus funges the route forward one step at a time and halts at the first finite step whose reading settles within the floor,
\[ N \;=\; \min\{\, n \in \mathbb{N} : d(r_n, r_{n-1}) \le \varepsilon \,\}, \qquad N < \infty. \]
The halt index $N$ is a definite finite place, not a limit. The apparatus checks the floor at each step and stops the first time the reading no longer moves by more than $\varepsilon$; it never inspects an infinite tail and never waits for a convergence it could not finish verifying. If no such $N$ arrives within the apparatus's budget, the route halts with a refusal rather than a promise of eventual settling.

The halt must not behave like impatience. A witness may want the route to end because the answer feels near, because the burden has grown tedious, or because a public report would be useful, but none of those wishes gives the route a valid stop. The apparatus halts only when the calibrated comparison has exposed a state that can answer or refuse under the route's own terms. Sometimes the halt records success, the carried receipt still answering after passage; sometimes refusal, the route broken at a local place; sometimes residue, the state answering one demand and leaving another open. The halt is a form of honesty, not a dramatic close: the route does not stop because nothing else could ever happen, but because the apparatus has enough return to expose what happened here. A later witness may reopen the claim under a new burden, but that reopening should start from a bounded state, not from a blur.

The halt also separates answer from exhaustion. Exhaustion would ask the route to spend every possible scene before anyone may speak, and a finite apparatus cannot wait for every possible burden before it reports. It must halt when the calibrated comparison has enough to expose a state under accountable terms. This does not make the answer absolute; it makes the answer usable. Later work may ask a new question, but it must do so against the halted report rather than against an unfinished haze. Halting is therefore the apparatus's first disciplined pause: for this route and this floor, here stands the answerable state --- a completion earned at a finite place, not a solution promised at an unreachable one.

The measured output names the third discipline. Once the route halts, the apparatus emits more than a verdict; a verdict alone says yes, no, or not-yet, while the measured output funges the halted state into a finite record that later hands can inspect. The output is the reading the route carried at the halt,
\[ \mathrm{out} \;=\; r_N \qquad (N < \infty), \]
a value emitted at a finite place, not a limit $r_\infty$ the apparatus could never reach. Around that value the report keeps the starting mark that anchored the claim, the extent of the passage, and the rate-like change the route recorded as it moved between states. These terms belong to a metaphysical scale rather than a physical trip: they record what the apparatus must preserve if another witness wants to inspect the route without swallowing every scene that produced it.

Each recorded term does definite work. The starting mark keeps the report tied to the claim that began the route, so a later witness can tell whether the route still answers the same demand. The extent records how far the route ranged through local scenes and burdens --- not an outer expanse, but how much accountable passage the claim survived. The rate-like change records how the state altered while the route moved; it is the finite response of the earlier section, the reading $\Delta_h$ taken as the route stepped, and it imports no outer travel-law. It helps the witness see whether the change stayed within the route's terms, crossed a refusal, or left residue another scene must face. Together with the floor, the halt condition, the answer, the refusal, and the residue, these give the halted route a finite surface. The floor must travel with the output, because $r_N$ answers only relative to the $\varepsilon$ that stopped it: a coarser floor halts sooner and emits a rougher value, a finer floor halts later and emits a sharper one, so a report that hid its floor would offer a number no later hand could weigh. The output is not an absolute reading but a reading good to a stated tolerance, and naming that tolerance is part of what makes the measurement public.

The measured output resists two failures at once. It must not shrink to a proclamation --- a bare claim that the route answered, asking acceptance without the route's terms --- and it must not swell into the whole passage, dragging every local detail forward until no later hand can carry it. The apparatus finds the middle by recording exactly the terms that make the answer repeatable in public: where the route began, what changed, what floor governed comparison, where the route halted, and what residue still has a name. That middle record also keeps the witness honest, because the starting mark and bounded terms remain available for another hand to challenge; a witness cannot quietly replace the claim with a more convenient one once the report exists.

The three disciplines now hold one another in place. Calibration gives the route a floor that later hands can revisit; accountable comparison orders the route's states against that floor; the halt closes the route at a finite place where extension would otherwise hide the answer; the measured output records the halted state in terms another witness can inspect. Together they let a traveling claim become a finite public answer without pretending that it has escaped locality. The claim still belongs to the scenes that answered or refused it, and the report does not erase those scenes but carries their answerable pattern forward as a bounded object.

This section closes before the next machinery begins. A calibrated route can halt; a halted route can emit; a measured output can preserve starting mark, extent, rate-like change, refusal, and residue. Yet the output has not explained why one original closing form should govern the apparatus, nor how later standing follows from that form: it has become inspectable, but not yet final. The next section asks what makes the apparatus's original close public and what standing a completed route can carry. A measurement is not a reading that has gone on long enough; it is a reading that has earned the right to stop, and the right to stop is the right to be answered for at a finite place a later hand can reach.

\section{Original Closure}

The bounded report from the previous section gives the route a finite surface.
It records the starting mark, the accountable extent, the rate-like change, the
governing baseline, the stop, the answer, the refusal, and the residue. In the chapter's own symbols that surface is a finite tuple,
\[ \mathrm{rep} \;=\; \big(\, c_0,\ \mathrm{ext},\ \Delta_h,\ \varepsilon,\ N,\ \mathrm{ans},\ \mathrm{ref},\ \mathrm{res} \,\big), \qquad N < \infty, \]
one field for each mark just named: the originating claim written $c_0$, the rate-like change carried as the finite response $\Delta_h$ of the earlier section, the governing floor $\varepsilon$, and the stop $N$ reached at a finite place. The
apparatus tanges the stopped route into that surface: it makes the route
holdable, so a second witness can take the route in hand without retracing it.
A holdable surface settles only half of what the route owes. The surface shows
that the apparatus had enough to stop. It does not yet show that the stop
belongs to the route as the route's own close rather than as a loose ending
pinned on from outside.

The remaining demand is a demand for return. A route earns its close when the
report comes back to the claim that began the route and answers from the carried
terms that gave the route public force. A stop fixes a finite place to receive
the account; closure adds the further discipline that the account, once
received, still answers the opening claim. Between the two lies the room where a
stopped report can drift: it may name a result while forgetting the mark that
first called for witness, keep a convenient summary while dropping the refusal
that made the summary honest, or preserve a residue as a word while losing the
scene that gave the residue its shape. Original closure is the work that holds
the report against that drift. The apparatus gathers the route without replacing
it, and hands a later witness a surface that still remembers what the route
claimed, where the claim entered, which baseline governed comparison, which
resistance the route met, and what stayed unresolved.

A report surface is what the apparatus hands over after stopping, and it carries
the public weight of a stopped route in finite form. A later witness cannot
receive the whole history at once. The next hand needs something small enough to
inspect, to challenge, to carry, and to pass onward, yet faithful enough that
the small thing still answers for the large route behind it. Tanging the route
is what makes such a surface available: a condensed face another witness can
hold.

Condensation is the delicate part, because a careless compression buys
smoothness by spending the route's honesty. A faithful surface keeps several
marks visible at once. It carries the starting claim, since standing belongs
only to the claim that actually entered; the receipt, since the claim reached
later witnesses through carriage rather than by assertion; the baseline, since
comparison without a public baseline decays into a memory of impressions; the
stop, since public use needs a finite place to receive the account; and the
refusal and residue together, since a route that survives with no remembered
risk has become ceremony. Held together, these marks let a later witness ask one
sharp question: what did the apparatus hand over, and how does that handover
answer the route that produced it?

The surface also keeps proportion. Some passages of the route carry more weight
than others. A refusal that exposed the claim's limit can matter more than three
easy returns. A small residue can govern how far the surface travels. A baseline
that looked ordinary at the start can become the only reason later comparison
still makes sense. So the apparatus arranges the surface by public force,
letting the heavier marks stand where a later witness will look first. A report
surface, in this sense, is less a list than an ordered face: the route made
holdable, with its costs left where the next witness can find them.

Closure also needs gathered completion. A completed route has gathered enough of
its own passage for one finite account to answer for the whole, even though
later work may still follow it. Earlier marks, comparisons, refusals, local
limits, and residues stop sitting as scattered fragments and take a shape that a
later witness can inspect. The gathering turns a sequence of encounters into a
single completed route, one that keeps its order, its costs, and its open edges
while wearing a face that can travel.

A well-gathered route holds plurality and unity together. Each scene keeps its
own answer, and the answers still belong to one carried claim. The local
differences that made resistance meaningful survive the gathering, while the
claim, receipt, refusal, and residue that traveled across the scenes give those
differences a common spine. That balance is what lets a later witness tell that
a single claim traveled, rather than that several unrelated performances merely
shared a name.

Completion belongs to the route, not to anyone's satisfaction with it. The
apparatus completes a route by gathering what the route itself carried and by
making that gathering answerable; an onlooker's sense that enough has been said
does not complete the route in the route's place. When a later witness asks why
the report counts, the apparatus answers through the gathered order: the claim
that entered, the receipt that carried it, the comparison that pressed it, the
refusal that could have stopped it, the residue that remained, and the return
that brings those marks back intact.

Original closure names that accountable return. A report closes originally when
it comes back to the starting claim without changing what the claim asked the
apparatus to carry. The close keeps the route's origin public: the first mark,
the carried receipt, the visible refusal, and the named residue all answer
together to the claim that opened the route. A close of this kind draws its
authority from the route's own terms, not from the route's length, the number of
witnesses it gathered, or the usefulness of the result. What it asserts is
strict and small: the route can now face its beginning again and show that the
report answers the claim that first entered, under the receipt and the refusals
that made the claim public. In the chapter's symbols, closure is the return of the report to that first claim,
\[ \mathrm{rep} \;\xrightarrow{\ \text{answers}\ }\; c_0 \qquad (N < \infty), \]
a finite relation rather than a limit: the report answers $c_0$ from the terms carried between, at the stop $N$, not at a fixpoint reached only beyond every stage.

The word \emph{original} points forward to a public origin, not backward into
private possession. The first mark never owned the later answer; it gave the
later answer something to answer to. Original closure returns to that mark as a
constraint. A route may cross several local scenes, narrow its promise, meet
refusals, and leave residue, and still close, provided the report shows how
those passages answer the first claim rather than escape it. The changes the
route underwent become the content of the close rather than a threat to it.

The carried receipt returns with the claim. A close that kept only the ending
would ask a later witness to trust that the route arrived while hiding how it
passed. Original closure keeps the receipt beside the claim, so the carried
order stays open to inspection: which terms survived, which narrowed, which met
refusal, and which still carry residue. The close leaves the route open to
challenge and gives that challenge a focused surface to grip.

Refusal returns as well. A report that hides its refusals makes success look
effortless and failure look accidental. Kept visible, a refusal tells a later
witness where the route risked itself. Where the route survived, the refusal
still names the pressure the claim answered. Where the route failed in part, the
refusal marks the boundary the report must carry. Where the route left residue,
the refusal helps locate the unfinished edge. Inside the close, refusal works as
a public guide to the cost of the answer.

Residue gives the close its honesty. A route can close without pretending that
every question has dissolved into a clean ending. The residue names what the
route did not settle, what a later burden may reopen, or what the report must
carry as a remainder. Here original closure parts from finality: finality would
end every question, while closure asks the route to return with enough
discipline that a later witness sees both the answer and the remainder. A report
keeps standing by keeping its residue in view, and loses that standing when it
buys smoothness by erasing its remainder.

So original closure is a bounded act. It does not reach for the whole world or
settle every later burden, and it makes no holder, crowd, or office into the
maker of truth. It closes this route by returning this report to this origin,
under this carried receipt, with this refusal and this residue still in view.
The narrowness is the strength. A report can travel precisely because the close
states exactly what it carries: another witness can receive it without
inheriting every hidden circumstance, and can challenge it without dissolving
the account into rumor.

Once original closure holds, the apparatus can set stages of the route beside
one another. An earlier stage and a later stage need not become the same event.
The first local answer may have met a single burden; the later report may gather
several pressures; the stopped surface may carry residue the first witness had
not yet named. Correspondence between stages is therefore a match of carried
relation rather than of event, scene, or wording: two stages correspond when the
same claim travels through different states while the relation among claim,
receipt, refusal, and residue stays accountable.

That carried relation says more than \emph{before} and \emph{after}. Sequence
alone gives order; the carried relation gives answerability across the order. The
earlier stage reports what the claim could carry at one place of return. The
later stage reports what the gathered route can carry after comparison, stop, and
report. The apparatus asks whether the two stages face one another under the same
public demand. When they do, the later report answers with the earlier stage
instead of replacing it. When they do not, the report names the break, the
narrowing, or the residue that keeps the stages from sharing one account.

Correspondence of this kind protects difference instead of erasing it. The
apparatus needs only the relation to carry, not the two scenes to play identical
roles. A receipt may answer in one scene as a first survival and in another as a
gathered close. A refusal may begin as a local wound and later serve as the
boundary of the report. A residue may appear first as a remainder and later as
the named edge that prevents overclaim. The stages differ while their relation
stays answerable, and that answerability lets the apparatus treat the route as
one route rather than a pile of disconnected moments.

A carried relation can also fail, and the apparatus gains by exposing the
failure rather than smoothing it. A later report may keep the same words while
changing the claim it pretends to answer, carry the receipt but drop the
refusal, or preserve the mark but hide the residue. Forced correspondence in
those cases would close a route that no longer returns to its own beginning. An
exposed failure keeps the close honest, since correspondence between stages
earns its force only when the later report can face the earlier stage without
disguise. The apparatus treats correspondence as something a route earns at each
return, not as a courtesy extended to any report that ends.

When the original close and the carried relation hold together, the report gains
inferred standing. Inferred standing does not float above the apparatus; it
follows inside the apparatus from the closed relation. The apparatus has a
stopped report, a gathered route, an accountable return, and stages that
correspond under one carried relation. From that order, the report earns
standing as a disciplined consequence, not from private certainty, a louder
voice, a perfect ending, or a final owner, but from the public way the route has
closed over its own carried terms.

The word \emph{inferred} stays modest. It means that standing follows from the
route's own accountable order, and nothing larger. It does not promise that
every possible question has an answer, that the apparatus now sees from beyond
locality, or that a later witness must keep silent. Standing follows because the
report returned to its origin, preserved its receipt, exposed refusal, carried
residue, and kept its relation across stages. A later witness may still challenge
that standing, but the challenge meets a closed relation rather than an unbounded
haze.

Standing is what lets the report funge. Before closure, the report stays
inspectable but not yet settled as the apparatus's own close. After closure, the
report funges: it circulates as a public consequence of the route, passing
between later witnesses without dragging the whole origin behind it. A later
participant receives the closed surface, inspects it, asks after the receipt,
presses the residue, and sets a new burden against the relation already carried.
The circulation is disciplined rather than free, since the report funges with its
costs attached, and what travels is the carried relation rather than a
smoothed-over version of it.

Inferred standing also spares the apparatus an endless restart. Without it, every
later witness would reopen each route from its first mark before any public
speech could begin, and shared work would never get past its own caution. With
bounded standing, a later witness begins from the closed report while keeping the
right to inspect and to refuse. The report travels as a disciplined consequence
rather than an unquestionable treasure; it gives public work a starting point
that already carries its costs, and it keeps later disagreement tied to the
carried relation rather than to a fog of unremembered scenes.

A later witness, then, never receives a bare answer. The witness receives a
closed surface with handles: the starting mark, the receipt, the baseline, the
stop, the refusal, the residue, and the carried relation. Each handle gives a
challenge an honest place to grip. A question about the mark sends the report
back to the claim that started the route. A question about the receipt asks what
traveled. A question about the stop asks why extension would have hidden rather
than clarified the answer. A question about the relation sets the earlier and
later stages side by side without pretending they became one event. Inferred
standing makes a challenge sharper rather than weaker, since it tells the
challenge where to begin.

The same standing lets the report funge beyond its first receiving scene. A
report that closes originally need not carry every local circumstance into every
later encounter; it carries the accountable surface those circumstances earned.
A new scene tests that surface against a fresh burden while respecting the
route's named terms. If the new burden presses the same carried relation, the
old standing gives the new route a disciplined place to start. If the burden
presses another edge, the report's residue shows the witness why a new route may
have to begin. Circulation follows a plain rule: carry the close while the
relation carries, and begin again when the relation no longer answers.

With this, the apparatus holds the original close that Chapter 5 needed. Truth
entered as witness rather than ownership. Reality and locality gave witness a
resistant place. Variation tested the claim without dissolving it. Paths and
mediation carried the answer across local scenes. Calibration set a public
baseline, comparison ordered the route's states, the stop fixed a finite place
to receive the account, and the bounded report made the stopped state holdable.
Original closure gathers that report back to the claim, receipt, refusal, and
residue that gave it force.

This closes the numbered sequence without closing the chapter's larger work. The
later essay, bridge, and coda can ask how the completed sequence should stand
within the book's broader argument, and can gather tone, transition, and
afterward pressure; the present section does not begin that work. It leaves the
apparatus with a public close: a report surface tanged from the stopped route, a
gathered completion, an accountable return, a correspondence across stages, and
an inferred standing that funges only as far as the apparatus that earned it
allows. The chapter can now turn from numbered construction to reflection on what
the construction has made visible.

\section*{Bridge: After Closure, the Shape of the Reading}

The chapter closed the apparatus on its own truth. The apparatus learned to witness a claim rather than own it, to test it where it stands, to carry it by mediation without losing what earned it, to calibrate a floor and halt its reading at a finite place, and to return its report to the claim that began it. Original closure was the last of these: the apparatus inferring its standing from its own earned record, the report answering the claim it opened with. That is a complete account of \emph{whether} the apparatus reads truth. It is not yet an account of the \emph{shape} of what the apparatus reads.

The distinction is the hinge to the next chapter. Closure settled that a reading can be witnessed, carried, halted, and returned to its claim; it did not ask what the reading looks like as the apparatus varies what it measures. A closed reading is a single finite value, the output $r_N$ the route emitted at its halt, and the apparatus knows that value is answerable. What it does not yet know is how that value behaves --- whether it holds steady, drifts, or turns sharply --- when the configuration it reads is moved a little to one side.

The apparatus is not unequipped for that question. In reading a claim locally it already learned to register the finite response of an answer to a nearby perturbation: the reading $\Delta_h$ taken at a definite step, not a limit pursued to vanishing. What it has not yet done is turn that instrument on its own settled output. When the apparatus tanges its closed reading and funges a controlled perturbation through it, it gets back not another verdict but a response, and the pattern of those responses across the configurations it can perturb is what the next chapter will call the shape of the reading.

The shift from \emph{whether} to \emph{shape} is a shift from a point to its neighborhood. A verdict answers at one configuration: this claim, read here, closed true. A shape answers across the configurations near it: how the reading would have closed had the apparatus measured slightly otherwise. The apparatus does not abandon the point --- it keeps the closed reading and its receipt --- but it now treats that point as the center of a small, surveyable region and asks what the reading does across it. The finite first variation is the first and coarsest record of that region: the leading way the reading bends as the configuration leaves the center.

The next chapter poses this as the question of stationarity. To ask when a reading is stationary is to ask when its finite first variation vanishes at a configuration --- when, to first order and within the floor that calibrated it, perturbing the configuration leaves the reading unmoved. This is what classical analysis would call the first derivative of the reading, kept here as the response $\Delta_h$ read at a definite step rather than a limit taken to zero. And the next chapter insists that stationary does not mean flat: a reading can be unchanging to first order while still carrying detail beneath that order, so that the vanishing of the first-order response is not the absence of content but a detector for it --- the apparatus learning to hear, in the silence of a first variation gone to zero, the presence of something the first order cannot show.

This is no idle refinement for a device whose work is to measure. A reading that lurches under the smallest perturbation answers for nothing beyond the single configuration it happened to meet, while a reading that holds steady across a neighborhood answers for the neighborhood as well. Stationarity is therefore the apparatus asking after the robustness of its own output: whether the value it closed on is an accident of where it stood or a stable feature of what it read. The shape of the reading, in this sense, measures how far a closed truth can be trusted to stay itself when the ground beneath it shifts.

This question follows from closure rather than departing from it. A reading the apparatus could not close would have no stable value to vary; only because the apparatus can halt at a finite place and return a definite output can it ask, of that output, how it changes. The closed apparatus therefore becomes, in the next chapter, an apparatus that reads shape, one that funges its settled reading through a controlled variation and watches what survives. The truth earned in this chapter is the precondition for the shape read in the next: a reading must be answerable before its steadiness can mean anything at all.

So the next chapter turns from closure to the first variation, opening the finite analytic behavior the apparatus's truth made possible: the transparent middle where a reading can be examined, the exact first variation it takes there, the balance that holds at that middle, and the single residual that will wear three names across the chapters ahead. None of that machinery belongs here; the bridge only marks the turn. The apparatus can now say that its readings are true, and it carries across this threshold the one question its own closure raises but cannot answer --- not whether a reading is true, but what a true reading is shaped like, and what holds steady in it when everything around it is allowed to move.

\section*{Coda: Witness in Review}

Chapter 5 earned a truth the apparatus could close on its own terms, and it earned it one rung at a time, each rung forced by a lack the rung before it exposed. Read forward, the chapter offers five capacities --- witness, locality, mediation, halting, and closure. Read in reverse, it shows that those capacities could not have come in any other order, and that nothing in the order was imported from outside the finite apparatus. This review walks that forced order once more, naming at each step the exact lack that made the next step unavoidable, and asking what the chapter has actually settled about what makes a reading true.

The chapter begins by refusing to let a claim be true because some position holds it. Truth is not a possession but a relation between the position that asserts a claim and a distinct position that receives it and answers for it. A claim held where it was made is a closed loop, $P \xrightarrow{\,c\,} P$, in which the asserter is the only receiver; such a loop has never been tested, because a position cannot witness its own claim. Witness is instead the directed passage to a different position that answers into a shared record,
\[ \text{witness:}\quad P \;\xrightarrow{\;c\;}\; Q \;\xrightarrow{\;\text{answers for } c\;}\; \text{record}, \qquad Q \neq P, \]
and the single inequality $Q \neq P$ carries the whole move: a truth that never leaves $P$ has not entered the public order at all. Yet witness, having made a claim public, says nothing about where the claim must answer. A witnessed claim can still float, dodging every challenge by shifting the place of its answer whenever pressure arrives. That missing place forces the next rung.

Locality supplies the place. A witnessed claim must answer from a bounded scene, so the apparatus tanges the claim at that scene and reads its finite response to a nearby perturbation,
\[ \Delta_h\, c(x) \;=\; \frac{c(x+h) - c(x)}{h}, \qquad h \neq 0, \]
the change in the claim's answer when one input is moved by a definite step $h$ --- a reading the apparatus can exhibit, sharpened along a chosen line of pressure to the directed response $D^h_v$, and taken always at a definite resolution rather than driven to a limit it could never finish taking. A claim that holds its shape under such pressure answers for its scene; a claim that breaks names the place where it broke. But a claim that has answered in one bounded scene has not shown that it can answer in another, and an answer left where it stood is parochial. That missing reach forces the next rung.

Mediation gives the local answer reach without dissolving it. A path is a finite chain of positions, and the mediation along it composes the local carries that funge the receipt from one position to the next, $m_\gamma = m_n \circ \cdots \circ m_1 : P_0 \to P_n$ with $n < \infty$, each carry preserving the conditions that gave the receipt its weight. What travels is the route back to the receipt, not the result alone: a distant position inherits the terms that earned the answer rather than a bare conclusion it would have to take on faith. The composition is finite --- the apparatus carries a truth, it does not re-assert one. But a route that can keep travelling yields no public result, since mediation can move an answer from scene to scene indefinitely without ever fixing a value a later hand can receive. That missing end forces the next rung.

Halting gives the route its end. The apparatus fixes a floor $\varepsilon$ below which two readings no longer count as disagreeing, so that two states meet when $d(r_m, r_n) \le \varepsilon$, and it funges the route forward one step at a time until the reading settles within that floor,
\[ N \;=\; \min\{\, n \in \mathbb{N} : d(r_n, r_{n-1}) \le \varepsilon \,\}, \qquad N < \infty, \]
halting at the first finite $N$ and emitting the value it carried there, $\mathrm{out} = r_N$ --- a reading good to a stated tolerance, not a limit $r_\infty$ the apparatus could never reach. The stop is completion, not interruption; the apparatus never inspects an infinite tail, and if no such $N$ arrives within its budget it halts with a refusal rather than a promise of eventual settling. But a halted, emitted reading is only inspectable so far; it has not yet shown that the stop belongs to the route as the route's own close rather than as a loose ending pinned on from outside. That missing return forces the last rung.

Original closure supplies the return. The apparatus tanges the route into a finite surface, the tuple $\mathrm{rep} = (\,c_0,\ \mathrm{ext},\ \Delta_h,\ \varepsilon,\ N,\ \mathrm{ans},\ \mathrm{ref},\ \mathrm{res}\,)$ with $N < \infty$, holding the originating claim, the extent, the finite response, the governing floor, the stop, the answer, the refusal, and the residue. A route closes originally when that surface returns to and answers the claim it began with,
\[ \mathrm{rep} \;\xrightarrow{\;\text{answers}\;}\; c_0 \qquad (N < \infty), \]
a finite relation rather than a limit: the report answers $c_0$ from the terms carried between, at the stop $N$, and draws its authority from the route's own record rather than from its length, its audience, or any holder's certainty. With this the apparatus closes on its own truth, and the chapter's forced order is complete.

Beneath every rung lies one object the chapter never multiplies and never replaces: the receipt the route carries. Witness enters it into a record a second position can read; locality tests it against a bounded scene; mediation carries it along a finite path; the halt fixes it at a finite place; and closure returns it to the claim it began with. No rung manufactures a second receipt to stand behind the first. The chapter grows by taking on further duties toward one carried passage, not by accumulating new passages each demanding its own ground, and this is why its truth can be strong and still finite: the strength is depth of obligation toward a single answerable record, not the length of a chain of fresh claims.

The order cannot be permuted. Closure presupposes a halted output to return, so it cannot precede the halt; the halt presupposes a route to arrest, so it cannot precede mediation; mediation presupposes a local answer to carry, so it cannot precede locality; locality presupposes a witnessed claim to place, so it cannot precede witness; and witness presupposes a claim distinct from the position that asserts it, so it rests on nothing earlier than itself. Each rung consumes the duty beneath it and produces the duty above. The five capacities are therefore not a sequence the chapter happened to prefer but the only sequence in which a finite apparatus can pass from a bare claim to a truth it can close --- a different order would not be a different route to the same end; it would be no route at all.

One feature runs through every rung: each is honestly finite. Witness needs two positions and a record, not a truth-predicate ranging over an infinite model. Locality reads a response at a definite step $h$, not a derivative taken in the limit. Mediation composes a path of length $n < \infty$, not a continuous transport through a space it could never traverse. Halting stops at a finite $N$, not at a convergence whose infinite tail it could never verify, and the output is the value $r_N$ rather than a limit $r_\infty$. The report is a finite tuple, not an unbounded archive. The apparatus never reaches for an infinite object to secure a finite truth, because a truth secured by a limit it could not exhibit would draw on an authority it could not redeem under challenge.

Read in reverse, the order also shows what it refuses to accept as truth. A claim does not become true because the apparatus can hold it, because a single position feels no doubt, because many positions repeat it, because it carries an old and respected name, because it ran a long time before stopping, or because it ended cleanly with no remainder. Each of those is a way of owning a claim rather than answering for it. Witness, locality, mediation, halting, and closure are precisely the disciplines that keep ownership from passing itself off as truth, and a report keeps its standing only while it holds its refusals and its residue in view beside its answer.

This is the turn on which the book pivots. The earlier chapters earned an apparatus that could record an answerable account --- a mark placed, a clock complemented, a representation bounded, a value carried under its overhead. Such an apparatus can say, with discipline, what it holds and what that holding cost. What it could not yet say is whether what it holds is true. Chapter 5 is where recording becomes truth, not by adding a faculty the apparatus lacked, but by demanding of the record it already kept that it answer a position other than the one that wrote it. Truth, on this account, is not a richer kind of record; it is the same record made to face a challenger and survive.

This fixes what the chapter means by a true reading, and what it does not. The apparatus does not certify a reading by matching it against a world no position can reach, nor by checking it for mere internal consistency, nor by counting the positions that happen to agree. A match to the inaccessible could never be exhibited; consistency alone would let a closed and self-pleasing account call itself true; and agreement is only ownership spread across many hands. What the chapter offers instead is answerability: a reading is true when it can be witnessed, tested where it stands, carried without loss, halted at a finite place, and returned to the claim it began with, so that any later position can take up the closed surface and press it again. Truth here is not a possession the apparatus secures once; it is a relation the apparatus keeps open.

The chapter's title named three things, and the forced order earns all three. Witness is truth received rather than owned, the first rung and the condition of the rest. Reality is what the bounded scene and the carried path made the claim answer to --- an order outside the apparatus's own preference, which can withhold the answer the claim wanted and so keep the reading from being a private decision. Decision is the calibrated halt and the original close, the two places where the apparatus must commit: that the reading has gone far enough to read, and that the report answers the claim it set out from. Witness opens the chapter, reality tests it, and decision closes it, and the three together are what separate a measurement that is merely recorded from one that is true.

So Chapter 5 leaves the apparatus able to close on its own truth, and to do so without owning it. What it earned is not a possession but a discipline: the standing of a claim follows from the public way the route closed over its own carried terms, and from nothing larger. A later position may still challenge that standing, but the challenge now meets a closed relation rather than an unbounded haze. Truth, in this apparatus, is what survives the distance between the position that asserts a claim and the position that answers for it. The apparatus did not learn to own that distance; it learned to keep it open, and to be answered across it.
